\scanpage{119}
% archive sheet number: 118
satisfaisant à \unsure{la|\deleted{la}|\deleted{la}la} condition que son carré $\kappa^{\otimes2}$ s'\deleted{$\unsure{\kappa^{\otimes2}}$}identifie, via les isom. canoniques
\[
 \underline{t}_E^{\otimes2}\simeq T_{\Sigma_{\unsure{I|J}},t},\qquad
 \mathcal O_t(1)^{\otimes2}=\mathcal O_t(2)\underset{\text{via }\alpha}{\simeq}
 \mathcal O_t(2)(\underline\omega)\simeq T_{\Sigma_I,t} \tag*{(77)}
\]
à l'isom identique de $T_{\Sigma_I,t}$ -- ce qui, pour $\alpha$ fixé (deux choix), détermine $\kappa$ modulo signe. Pour $\alpha$ non fixé, $\kappa$ est déterminé modulo multiplication par $\mu_4(k)$.

Rappelons-nous de la suite exacte canonique
\diagram{\begin{tikzcd}[ampersand replacement=\&,column sep=small,row sep=small]
 (78) \& 0\arrow[r] \& F_t\arrow[r]\arrow[d,phantom,"\simeq"{description,rotate=90}] \& V_I(k)\arrow[r]\arrow[d,phantom,"\simeq"{description,rotate=90}] \& \mathcal O_t(1)\arrow[r] \& 0\\
 \& \& \mathcal O_t(-1)(\underline\omega) \& \Gamma(\Sigma_I,\underline{\mathcal O}_{\Sigma_I}(1))
\end{tikzcd}}
que \deleted{\unsure{par tenso}} je compare à
\diagram{\begin{tikzcd}[ampersand replacement=\&,column sep=small]
 (79) \& 0\arrow[r] \& k.1_{\Sigma_I}\arrow[r] \&
 \overbrace{\Gamma(\Sigma_I,\underbrace{\underline{\mathcal O}_{\Sigma_I}\{t\}}_{\substack{\simeq\underline{\mathcal O}_{\Sigma_I}(1)\\\text{non canoniquement}}})}^{\simeq\unsure{\mathcal F_2|F_2}}\arrow[r] \&
 T_{\Sigma_I,t}\arrow[r] \& 0
\end{tikzcd}}
On constate que la deuxième suite exacte est canoniquement isom. à celle déduite de la première, en la tensorisant par
\[
 \check F_t\simeq\mathcal O_t(1)(\underline\omega) \tag*{(\deleted{\unsure{79}}80)}
\]
(on a un iso valable chaque fois qu'on a un Module loc libre de rang $2$, $\underline V$, \deleted{\unsure{dont on|don}} \inserted{et} une section $t$ de $\mathbb P(\underline V)$ \ldots) :
\diagram{\begin{tikzcd}[ampersand replacement=\&,column sep=small,row sep=normal]
 \deleted{(\unsure{79|80})}(81) \& 0\arrow[r,dash] \& F_t\otimes\check F_t\arrow[r]\arrow[d,"\sim"] \& V_I(k)\otimes\check F_t\arrow[r]\arrow[d,"\sim"] \& \overbrace{\mathcal O_t(1)\otimes\check F_t}^{\underline{\mathcal O}_t(2)(\underline\omega)}\arrow[r]\arrow[d,"\sim"] \& 0\\
 \&0\arrow[r]\& k.1_{\Sigma_I}\arrow[r] \& \underbrace{\Gamma(\Sigma_I,\underline{\mathcal O}_{\Sigma_I}\{t\})}_{\unsure{\mathcal F_2}}\arrow[r]\&T_{\Sigma_I,t}\arrow[r]\&0.
\end{tikzcd}}
